\textbf{Phần 1. Trắc nghiệm lựa chọn}
\setcounterpageref{ex}{0}

\begin{ex} %[2D2N08]
Trong bảng cộng (qua 10), phép tính $9 + 9$ có kết quả là
\choice
{$17$}
{\True $18$}
{$19$}
{$20$}
\loigiai{
$9 + 9 = 18$.\\
\textbf{Đáp án của câu hỏi là B}
}
\end{ex}

\begin{ex} %[2D2N08]
Phép tính nào sau đây có kết quả bằng 11?
\choice
{$9 + 3$}
{$8 + 4$}
{$7 + 5$}
{\True $9 + 2$}
\loigiai{
$9 + 2 = 11$.\\
\textbf{Đáp án của câu hỏi là D}
}
\end{ex}

\begin{ex} %[2D2H08]
Phép tính nào sau đây có kết quả bằng 13?
\choice
{$8 + 4$}
{\True $8 + 5$}
{$9 + 5$}
{$7 + 7$}
\loigiai{
$8 + 5 = 13$.\\
\textbf{Đáp án của câu hỏi là B}
}
\end{ex}

\begin{ex} %[2D2H08]
Phép tính nào sau đây có kết quả lớn nhất?
\choice
{$8 + 7$}
{$9 + 7$}
{$8 + 8$}
{\True $9 + 8$}
\loigiai{
$9 + 8 = 17$ (lớn nhất).\\
\textbf{Đáp án của câu hỏi là D}
}
\end{ex}

\begin{ex} %[2D2H08]
Phép tính nào sau đây có kết quả bé nhất?
\choice
{$8 + 4$}
{$7 + 5$}
{\True $6 + 5$}
{$9 + 3$}
\loigiai{
$6 + 5 = 11$ (bé nhất).\\
\textbf{Đáp án của câu hỏi là C}
}
\end{ex}

\begin{ex} %[2D2H08]
Điền số thích hợp: $7 + \dots = 16$.
\choice
{\True $9$}
{$8$}
{$7$}
{$6$}
\loigiai{
$7 + 9 = 16$.\\
\textbf{Đáp án của câu hỏi là A}
}
\end{ex}

\begin{ex} %[2D2V08]
Có bao nhiêu phép cộng hai số có 1 chữ số có kết quả bằng 12? (Ví dụ: $9+3, 8+4, 7+5, 6+6, 5+7, 4+8, 3+9$)
\choice
{$5$ phép tính}
{$6$ phép tính}
{\True $7$ phép tính}
{$8$ phép tính}
\loigiai{
Các phép tính: $9+3, 8+4, 7+5, 6+6, 5+7, 4+8, 3+9$ (có 7 phép tính).\\
\textbf{Đáp án của câu hỏi là C}
}
\end{ex}

\begin{ex} %[2D2V08]
Tính: $8 + 7 + 3 = \dots$
\choice
{$17$}
{$19$}
{$20$}
{\True $18$}
\loigiai{
$8 + (7 + 3) = 8 + 10 = 18$.\\
\textbf{Đáp án của câu hỏi là D}
}
\end{ex}

\begin{ex} %[2D2V08]
Điền dấu thích hợp: $7 + 8 \dots 9 + 6$.
\choice
{$>$}
{\True $=$}
{$<$}
{Không so sánh được}
\loigiai{
$7 + 8 = 15$ và $9 + 6 = 15$. Vậy $7 + 8 = 9 + 6$.\\
\textbf{Đáp án của câu hỏi là B}
}
\end{ex}

\begin{ex} %[2D2V08]
Một luống hoa có 8 cây hoa hồng và 6 cây hoa cúc. Số cây hoa có tất cả là
\choice
{\True $14$ cây}
{$13$ cây}
{$15$ cây}
{$16$ cây}
\loigiai{
$8 + 6 = 14$ cây.\\
\textbf{Đáp án của câu hỏi là A}
}
\end{ex}

\begin{ex} %[2D2V08]
Tìm $x$ biết: $x - 7 = 9$.
\choice
{$15$}
{$17$}
{\True $16$}
{$14$}
\loigiai{
$x = 9 + 7 = 16$.\\
\textbf{Đáp án của câu hỏi là C}
}
\end{ex}

\begin{ex} %[2D2C08]
Tổng của hai số có 1 chữ số là 17. Hai số đó là
\choice
{\True $9$ và $8$}
{$8$ và $8$}
{$9$ và $9$}
{$7$ và $9$}
\loigiai{
Vì các số có 1 chữ số nên $9 + 8 = 17$.\\
\textbf{Đáp án của câu hỏi là A}
}
\end{ex}

\textbf{Phần 2. Trắc nghiệm đúng sai}
\setcounterpageref{ex}{0}

\begin{ex} %[2D2H08]
Xét tính đúng sai của các phép tính trong bảng cộng:
\choiceTF
{\True $9 + 4 = 13$.}
{\True $8 + 6 = 14$.}
{$7 + 8 = 16$.}
{\True $7 + 8 = 15$.}
\loigiai{
\begin{itemize}
    \item $9 + 4 = 13$: Đúng.
    \item $8 + 6 = 14$: Đúng.
    \item $7 + 8 = 15 \neq 16$. Ý 3 Sai, Ý 4 Đúng.
\end{itemize}
\begin{itemize}
    \item a) \textbf{Đúng}.
    \item b) \textbf{Đúng}.
    \item c) \textbf{Sai}.
    \item d) \textbf{Đúng}.
\end{itemize}
}
\end{ex}

\begin{ex} %[2D2V08]
Cho các phép tính: (1) $9 + 5$, (2) $8 + 6$, (3) $7 + 7$. Xét tính đúng sai:
\choiceTF
{\True Phép tính (1) có kết quả bằng 14.}
{\True Phép tính (2) có kết quả bằng 14.}
{Phép tính (3) có kết quả lớn hơn phép tính (1).}
{\True Cả 3 phép tính đều có kết quả bằng 14.}
\loigiai{
\begin{itemize}
    \item $9 + 5 = 14$: Đúng.
    \item $8 + 6 = 14$: Đúng.
    \item Cả 3 đều bằng 14. Ý 3 Sai, Ý 4 Đúng.
\end{itemize}
\begin{itemize}
    \item a) \textbf{Đúng}.
    \item b) \textbf{Đúng}.
    \item c) \textbf{Sai}.
    \item d) \textbf{Đúng}.
\end{itemize}
}
\end{ex}

\textbf{Phần 3. Trả lời ngắn}
\setcounterpageref{ex}{0}

\begin{ex} %[2D2H08]
Tính kết quả: $6 + 9$.
\shortans{15}
\loigiai{
\textbf{Đáp số:} $15$\\
$6 + 9 = 15$.
}
\end{ex}

\begin{ex} %[2D2H08]
Tính kết quả: $8 + 9$.
\shortans{17}
\loigiai{
\textbf{Đáp số:} $17$\\
$8 + 9 = 17$.
}
\end{ex}

\begin{ex} %[2D2V08]
Tính giá trị biểu thức: $9 + 5 + 1$.
\shortans{15}
\loigiai{
\textbf{Đáp số:} $15$\\
$(9 + 1) + 5 = 15$.
}
\end{ex}

\begin{ex} %[2D2V08]
Điền số thích hợp: $8 + ? = 17$.
\shortans{9}
\loigiai{
\textbf{Đáp số:} $9$\\
$17 - 8 = 9$.
}
\end{ex}

\begin{ex} %[2D2V08]
Bình có 7 cái nhãn vở, mẹ cho Bình thêm 8 cái nhãn vở. Hỏi Bình có tất cả bao nhiêu cái nhãn vở?
\shortans{15}
\loigiai{
\textbf{Đáp số:} $15$\\
$7 + 8 = 15$ cái.
}
\end{ex}

\begin{ex} %[2D2C08]
Tìm tổng lớn nhất có thể của hai số có 1 chữ số.
\shortans{18}
\loigiai{
\textbf{Đáp số:} $18$\\
$9 + 9 = 18$.
}
\end{ex}
